%% fig-system.tex -- Figure: high-level system architecture and component
%% interaction. One ASGI process, the clients that reach it, and the systems it
%% integrates with but does not own.
%%
%% Routing convention: every long edge runs either in the gutter between the
%% process band and the state column (lane A), in the outer lane to the right of
%% everything (lane B), or in the reserved channel below the band. Nothing
%% crosses a box.
%%
%% Type scale: \figH headings, \figB detail (7pt floor). Row detail lines were cut
%% to a single short clause each to pay for the larger type; the full enumerations
%% live in the prose of Sections 5 and 6.

\begin{figure}[tbp]
\centering
\ifcbexcalidraw
  \begin{cbwide}\cbdiagram{fig-system}\end{cbwide}
\else
\begin{cbwide}
\begin{tikzpicture}[x=1cm,y=1cm]

  \def\MX{3.30}   % process band left edge
  \def\MW{9.90}   % process band width
  \def\RX{13.66}  % state column left edge
  \def\RW{4.06}   % state column width
  \def\GA{13.36}  % lane A: control plane out to MLflow
  \def\GB{17.70}  % lane B: the production agent's round trip
  \def\CH{-7.72}  % channel below the band for long horizontal edges

  % ======================= clients ==========================================
  \node[anchor=north west, font=\figB\bfseries, text=cbMuted] at (0,0.46) {CLIENTS};
  \node[surface, anchor=north west, text width=2.44cm, minimum height=0.68cm]
    (br) at (0,-0.10) {\textbf{Browser}\\{\figB React SPA session}};
  \node[surface, anchor=north west, text width=2.44cm, minimum height=0.68cm]
    (ap) at (0,-0.98) {\textbf{API client}\\{\figB account-scoped}};
  \node[surface, anchor=north west, text width=2.44cm, minimum height=0.68cm]
    (sv) at (0,-1.86) {\textbf{Service caller}\\{\figB published workflow}};
  \node[extern, anchor=north west, text width=2.44cm, minimum height=0.92cm]
    (pa) at (0,-2.84)
    {\textbf{Production agent}\\{\figB the consumer, outside\\the control plane}};

  % ======================= the ASGI process =================================
  \node[band, anchor=north west, minimum width=\MW cm, minimum height=7.34cm]
    (proc) at (\MX,0.16) {};
  \node[bandlab, anchor=north west, text width=9.6cm, align=left] at (\MX+0.12,0.50)
    {ONE \sysbare{} ASGI APPLICATION \normalfont\itshape (embedded in \mlflow{}
     or standalone --- \figref{fig:topologies})};

  % ---- admission -----------------------------------------------------------
  \node[govern, anchor=north west, text width=\MW cm-0.60cm, minimum height=0.56cm]
    (adm) at (\MX+0.24,-0.10)
    {\textbf{Admission}\ \ {\figB session or token \dt scopes \dt CSRF \dt rate
     limit}};

  % ---- surfaces ------------------------------------------------------------
  \foreach [count=\i from 0] \t in {%
      {\textbf{Route handlers}\\{\figB \statRouteDecls{} declarations}},
      {\textbf{SSE + webhooks}\\{\figB live events}},
      {\textbf{\aria tool loop}\\{\figB risk-tiered}}}
  {
    \pgfmathsetmacro{\w}{(\MW-0.48-2*0.18)/3}
    \node[surface, anchor=north west, text width=\w cm-8pt, minimum height=0.78cm,
          inner xsep=2pt] at ({\MX+0.24+\i*(\w+0.18)}, -0.80) {\t};
  }

  % ---- lifecycle modes -----------------------------------------------------
  \node[control, anchor=north west, text width=\MW cm-0.60cm, minimum height=0.50cm,
        font=\figB] (modes) at (\MX+0.24,-1.72)
    {\mode{Author}\arrowto\mode{Test}\arrowto\mode{Evaluate}\arrowto%
     \mode{Calibrate}\arrowto\mode{Release}\arrowto\mode{Observe}};

  % ---- control-plane services ----------------------------------------------
  \foreach [count=\i from 0] \t in {%
      {\textbf{Orchestrator}\\{\figB triage, diagnose,\\candidate}},
      {\textbf{Evaluation}\\{\figB scorecards,\\LLM judges}},
      {\textbf{Calibration}\\{\figB \statOptimizers{} optimizer\\paths, replay}},
      {\textbf{Apply +}\\\textbf{promote}\\{\figB gate, alias,\\checkpoint}}}
  {
    \pgfmathsetmacro{\w}{(\MW-0.48-3*0.18)/4}
    \node[control, anchor=north west, text width=\w cm-8pt, minimum height=1.30cm,
          inner xsep=2pt] at ({\MX+0.24+\i*(\w+0.18)}, -2.36) {\t};
  }

  % ---- governance substrate ------------------------------------------------
  \node[govern, anchor=north west, text width=\MW cm-0.60cm, minimum height=0.48cm,
        font=\figB] (gov) at (\MX+0.24,-3.82)
    {execution policy \dt gate verdicts \dt release control \dt audit log \dt
     effect ledger};

  % ---- kernel --------------------------------------------------------------
  \node[control, anchor=north west, text width=\MW cm-0.60cm, minimum height=0.48cm,
        font=\figB] (kern) at (\MX+0.24,-4.44)
    {config \& secrets \dt persistence \dt storage \dt tool execution \dt event
     bus \dt provider adapters};

  % ---- durable queues + loops ---------------------------------------------
  \node[db, anchor=north west, text width=2.20cm, minimum height=1.86cm]
    (q) at (\MX+0.24,-5.10)
    {\textbf{Durable}\\\textbf{queues}\\[0.4ex]{\figB status, claim\\and lease
      columns\\on rows}};
  \node[async, anchor=north west, text width=\MW cm-3.60cm, minimum height=1.86cm]
    (loops) at (\MX+0.24+2.86,-5.10)
    {\textbf{Up to \statLoops{} in-process loops}\\{\figB no separate worker tier}\\[0.4ex]
     {\figB Refinement \dt CalibrationDrain \dt WorkflowRun\\
      AriaPlan \dt KnowledgeBase \dt Scheduler\\
      Janitor \dt WebhookDispatcher}};
  \draw[flow, draw=cbAsync] ($(q.east)+(0,0.24)$) -- ($(loops.west)+(0,0.24)$);
  \draw[flow, draw=cbAsync] ($(loops.west)-(0,0.24)$) -- ($(q.east)-(0,0.24)$);

  % ======================= state and external systems =======================
  \node[anchor=north west, font=\figB\bfseries, text=cbMuted] at (\RX,0.46)
    {STATE \& EXTERNAL};
  \node[extern, anchor=north west, text width=\RW cm-8pt, minimum height=0.94cm]
    (mlf) at (\RX,-0.10)
    {\textbf{\mlflow{} 3.14+}\\{\figB prompt versions + aliases\\traces \dt
     assessments}};
  \node[db, anchor=north west, text width=\RW cm-8pt, minimum height=0.94cm]
    (pg) at (\RX,-1.24)
    {\textbf{PostgreSQL 17}\\{\figB \statModels{} tables --- metadata\\\emph{and}
     the file inventory}};
  \node[db, anchor=north west, text width=\RW cm-8pt, minimum height=0.76cm]
    (obj) at (\RX,-2.38)
    {\textbf{Object storage}\\{\figB file bytes only}};
  \node[extern, anchor=north west, text width=\RW cm-8pt, minimum height=0.76cm]
    (llm) at (\RX,-3.34)
    {\textbf{LLM providers}\\{\figB OpenAI \dt Claude \dt Gateway}};
  \node[extern, anchor=north west, text width=\RW cm-8pt, minimum height=0.76cm]
    (mcp) at (\RX,-4.30)
    {\textbf{MCP servers}\\{\figB containment or sidecar}};
  \node[extern, anchor=north west, text width=\RW cm-8pt, minimum height=0.76cm]
    (tx) at (\RX,-5.26)
    {\textbf{Event transport}\\{\figB NATS \dt Redis \dt DB}};

  % ======================= edges ============================================
  % Account clients and service callers all funnel through admission.
  \draw[flow] (br.east) -- ({\MX+0.22},-0.32);
  \draw[flow] (ap.east) -- ({\MX+0.22},-0.38);
  \draw[flow] (sv.east) -- ({\MX+0.22},-0.44);

  % The layer spine: each row hands off to the one below it.
  \foreach \ya/\yb in {-0.66/-0.76, -1.58/-1.68, -2.22/-2.32,
                       -3.66/-3.78, -4.30/-4.40, -4.92/-5.06}
    { \draw[flow, draw=cbMuted] ({\MX+\MW/2},\ya) -- ({\MX+\MW/2},\yb); }

  % Only the kernel touches state; lane A carries the MLflow-owned effects.
  \draw[flow] (kern.east) -- (pg.west);
  \draw[flow] (kern.east) -- (obj.west);
  \draw[flow] (kern.east) -- (llm.west);
  \draw[flow] (kern.east) -- (mcp.west);
  \draw[flow] (loops.east) -- (tx.west);
  \draw[flow] ({\MX+\MW-0.12},-2.90) -- (\GA,-2.90) -- (\GA,-0.44) -- (mlf.west);

  % Lane B: the round trip that closes the loop, kept outside both the process
  % band and the state column so that it crosses nothing.
  \draw[{Stealth[length=4pt,width=3pt]}-{Stealth[length=4pt,width=3pt]},
        draw=cbExtern, line width=0.8pt]
    (pa.south) -- (1.22,\CH) -- (\GB,\CH) -- (\GB,-0.60) -- (mlf.east);
  \node[anchor=south, font=\figB, text=cbInk, align=center, inner sep=1.8pt]
    at (9.6,\CH+0.04)
    {the agent resolves \prodalias{} at call time and emits traces --- when
     \mode{Release} rotates the alias, the next call runs the new version with
     \textbf{no client redeploy}};

\end{tikzpicture}
\end{cbwide}
\fi
\caption{High-level architecture and component interaction. Every account client
and service caller enters through one admission stage; \Lmode{} drives four
control-plane service groups; and the kernel is the only tier that touches state
or external systems. Two properties are load-bearing for the rest of the paper.
First, \sys is a \emph{sibling} surface to \mlflow{}, not a proxy in front of it:
each owns its own store, which is what makes the reconciliation boundary in
\figref{fig:dataflow} necessary. Second, there is no separate worker tier ---
queued work is arbitrated through claim and lease columns on relational rows, so
one process both serves requests and drains queues. The paired arrows between the
queue and the loop group are a consumer's atomic claim and the status transition it
writes back (\algref{alg:claim}). The
outer lane is the only path that leaves the platform entirely, and it is why a
release needs no client deployment.}
\label{fig:system}
\figkey
\end{figure}
